\phantomsection
\addcontentsline{toc}{chapter}{原课程信息与作业安排}
\par\Needspace{4\baselineskip}
\noindent\textbf{原课程信息与作业安排}\par
MIT 18.650 \textit{Statistics for Applications}, Fall 2016, 是数学系本科课程. 原课程教师为 Philippe Rigollet; 初讲列出的助教为 Victor-Emmanuel Brunel. 课程从实际数据问题建立统计模型, 研究方法的数学依据、适用条件和局限. 本书前十章保留已整理的讲义主线, 后十一章依原次序编入 PS1--PS11 的完整中文题面与解答.

\Needspace{5\baselineskip}
\begin{longtable}{@{}p{3cm}p{13.8cm}@{}}
\toprule
项目 & 原课程公开信息\\
\midrule
网页大纲 & 每周两次课, 每次 1.5 小时. 先修为 18.440 \textit{Probability and Random Variables} 水平的概率论, 以及矩阵、向量和特征值等线性代数知识.\\
初讲的先修 & 概率课 18.600 或 6.041, Calculus 2, 以及矩阵、向量、乘法和正交等线性代数基础. 原片说明没有指定必读教材.\\
初讲的时段 & 周二、周四, 原片写为 \texttt{1:00--2:30am}. \texttt{am} 确实见于原片; 这里保留其标注, 不另推定时区或改作 \texttt{pm}. 选修习题课时间待定.\\
网页评分比例 & 作业 20\%, 期中考试 30\%, 期末考试 50\%.\\
初讲评分比例 & 作业 30\%, 期中考试 30\%, 期末考试 40\%. 每周作业一次, 共 11 次, 取最好 10 次.\\
初讲的考试安排 & 期中考试列为 11 月 8 日, 随堂, 80 分钟, 闭卷且不可带笔记, 可带 \textit{cheatsheet}. 期末时间待定, 2 小时, 可带书和笔记.\\
\bottomrule
\end{longtable}
网页大纲与初讲原片的评分比例、先修课程号存在差异, 上表分别列出, 不把其中一版静默替换为另一版. 具体定位为\href{https://ocw.mit.edu/courses/18-650-statistics-for-applications-fall-2016/pages/syllabus/}{官方 Syllabus}, 以及 \textit{Introduction} 原 PDF 的物理第 3--5 页. 原课确有期中和期末考核安排; 本次所核公开课程页面未发现考试题卷 PDF, 因而本书不虚构其题面或答案.

\par\Needspace{4\baselineskip}
\noindent\textbf{主题与讲次}\par
\href{https://ocw.mit.edu/courses/18-650-statistics-for-applications-fall-2016/pages/lecture-slides/}{原讲义内容页}按下列讲次分组. 公开课程导航未提供带逐次日期的独立课历; 讲次次序与作业截止日期须分开理解.
\begin{longtable}{@{}p{2cm}p{14.8cm}@{}}
\toprule
讲次 & 原课主题\\
\midrule
1--2 & 统计引论\\
3 & 参数推断\\
4--5 & 极大似然估计\\
6 & 矩估计\\
7--10 & 参数假设检验\\
11--12 & 拟合优度检验\\
13--16 & 回归\\
17--18 & 贝叶斯统计\\
19--20 & 主成分分析\\
21--24 & 广义线性模型\\
\bottomrule
\end{longtable}

\par\Needspace{4\baselineskip}
\noindent\textbf{作业日期与解答身份}\par
\href{https://ocw.mit.edu/courses/18-650-statistics-for-applications-fall-2016/pages/assignments/}{官方 Assignments} 列出 11 份作业, 并明确没有公开官方解答. 以下题面译自实际原件, 保留主问题、各层子问、无编号要求及原题勘误提示. 解答由 AI 独立推导并经本项目复审, 不是 MIT 发布的答案. 原题条件不足或有印刷误式时, 题面与整理注分开, 解答明确说明可证范围和修订版本.

各份原件均印为 Fall 2016, 只列截止日期, 没有单独的发布日或明确时区. 下表的日期依原件, 每次均为周五中午 12 点, 不把它们改写为其他学期或授课日期.
\begin{longtable}{@{}p{2cm}p{6.5cm}p{8.2cm}@{}}
\toprule
作业 & 原件截止日期 & 主问题主题\\
\midrule
PS1 & 2016 年 9 月 16 日, 中午 12 点 & 随机变量收敛、统计断言与 Bernoulli 置信区间\\
PS2 & 2016 年 9 月 23 日, 中午 12 点 & 偏差、可辨识性、Poisson 与均匀分布的置信区间\\
PS3 & 2016 年 9 月 30 日, 中午 12 点 & 极大似然估计及一致性、KL 与全变差距离\\
PS4 & 2016 年 10 月 7 日, 中午 12 点 & 极大似然、Fisher 信息、矩估计与删失数据\\
PS5 & 2016 年 10 月 14 日, 中午 12 点 & 检验与置信区间、双样本均值、隐含的检验判断\\
PS6 & 2016 年 10 月 21 日, 中午 12 点 & 检验基础、单侧 Student 检验、Bernoulli 独立性\\
PS7 & 2016 年 10 月 28 日, 中午 12 点 & QQ 图、双样本 KS 检验、连续分布的独立性\\
PS8 & 2016 年 11 月 4 日, 中午 12 点 & 异方差回归、随机设计回归与 logistic 回归\\
PS9 & 2016 年 11 月 18 日, 中午 12 点 & 固定设计的非参数回归与密度估计\\
PS10 & 2016 年 12 月 2 日, 中午 12 点 & 贝叶斯估计、贝叶斯线性回归与协方差矩阵\\
PS11 & 2016 年 12 月 9 日, 中午 12 点 & 指数族、单参数典则指数族与隐变量线性模型\\
\bottomrule
\end{longtable}
